%% Generated by Sphinx.
\def\sphinxdocclass{report}
\documentclass[letterpaper,10pt,english]{sphinxmanual}
\ifdefined\pdfpxdimen
   \let\sphinxpxdimen\pdfpxdimen\else\newdimen\sphinxpxdimen
\fi \sphinxpxdimen=.75bp\relax
\ifdefined\pdfimageresolution
    \pdfimageresolution= \numexpr \dimexpr1in\relax/\sphinxpxdimen\relax
\fi
\newdimen\sphinxremdimen\sphinxremdimen = 10pt
%% let collapsible pdf bookmarks panel have high depth per default
\PassOptionsToPackage{bookmarksdepth=5}{hyperref}

\PassOptionsToPackage{booktabs}{sphinx}
\PassOptionsToPackage{colorrows}{sphinx}

\PassOptionsToPackage{warn}{textcomp}
\usepackage[utf8]{inputenc}
\ifdefined\DeclareUnicodeCharacter
% support both utf8 and utf8x syntaxes
  \ifdefined\DeclareUnicodeCharacterAsOptional
    \def\sphinxDUC#1{\DeclareUnicodeCharacter{"#1}}
  \else
    \let\sphinxDUC\DeclareUnicodeCharacter
  \fi
  \sphinxDUC{00A0}{\nobreakspace}
  \sphinxDUC{2500}{\sphinxunichar{2500}}
  \sphinxDUC{2502}{\sphinxunichar{2502}}
  \sphinxDUC{2514}{\sphinxunichar{2514}}
  \sphinxDUC{251C}{\sphinxunichar{251C}}
  \sphinxDUC{2572}{\textbackslash}
\fi
\usepackage{cmap}
\usepackage[T1]{fontenc}
\usepackage{amsmath,amssymb,amstext}
\usepackage{babel}



\usepackage{tgtermes}
\usepackage{tgheros}
\renewcommand{\ttdefault}{txtt}



\usepackage[Bjarne]{fncychap}
\usepackage{sphinx}

\fvset{fontsize=auto}
\usepackage{geometry}


% Include hyperref last.
\usepackage{hyperref}
% Fix anchor placement for figures with captions.
\usepackage{hypcap}% it must be loaded after hyperref.
% Set up styles of URL: it should be placed after hyperref.
\urlstyle{same}

\addto\captionsenglish{\renewcommand{\contentsname}{Contents:}}

\usepackage{sphinxmessages}
\setcounter{tocdepth}{1}



\title{CS4100 Final Project}
\date{Apr 06, 2026}
\release{0.1.0}
\author{Christian Garcia, Tommaso Maga, Yu\sphinxhyphen{}Chun Ou, Peter SantaLucia}
\newcommand{\sphinxlogo}{\vbox{}}
\renewcommand{\releasename}{Release}
\makeindex
\begin{document}

\ifdefined\shorthandoff
  \ifnum\catcode`\=\string=\active\shorthandoff{=}\fi
  \ifnum\catcode`\"=\active\shorthandoff{"}\fi
\fi

\pagestyle{empty}
\sphinxmaketitle
\pagestyle{plain}
\sphinxtableofcontents
\pagestyle{normal}
\phantomsection\label{\detokenize{index::doc}}



\chapter{cs4100\sphinxhyphen{}pcr\sphinxhyphen{}gradcam\sphinxhyphen{}prediction}
\label{\detokenize{index:cs4100-pcr-gradcam-prediction}}
\sphinxAtStartPar
\sphinxhref{https://creativecommons.org/licenses/by-nc/4.0/}{\sphinxincludegraphics{{e296f812318df3c92058961183e28f85a819118b}.png}}

\sphinxAtStartPar
\sphinxstylestrong{Contributors:} Christian Garcia, Tommaso Maga, Yu\sphinxhyphen{}Chun Ou, Peter SantaLucia


\section{License}
\label{\detokenize{index:license}}
\sphinxAtStartPar
This project is licensed under \sphinxhref{https://creativecommons.org/licenses/by-nc/4.0/}{CC BY\sphinxhyphen{}NC 4.0} to comply with the licensing terms of the BreastDCEDL dataset, which itself adopts CC BY\sphinxhyphen{}NC 4.0 as the most restrictive license among its three source datasets from TCIA. Specifically, while I\sphinxhyphen{}SPY 1 (CC BY 3.0) and I\sphinxhyphen{}SPY 2 (CC BY 4.0) permit commercial use, the Duke Breast Cancer MRI dataset (CC BY\sphinxhyphen{}NC 4.0) does not. As a derivative work integrating all three sources, this project inherits that restriction. Academic use, research, and adaptation with attribution are welcome.


\section{Purpose}
\label{\detokenize{index:purpose}}
\sphinxAtStartPar
The group has decided to predict patient breast cancer outcomes using MRI scans and patient metadata. An accurate measure of these outcomes is given by the pathological complete response (PCR) metric for triple\sphinxhyphen{}negative breast cancer (TNBC). That is, the complete remission of an invasive cancer that is found in a tissue sample. For TNBC Patients with PCR, 90\% of patients experience event\sphinxhyphen{}free survival (EFS) over a period of three years, while only 67\% experienced EFS over the same period (Toss, et al.). By predicting PCR for some patients, we are able to accurately predict patient outcomes over a period of three years.
By having an accurate prediction of a patient’s PCR, doctors can preemptively consider other treatment plans/evaluate the need for surgery for some patient. Given the aggressiveness of TNBC, it’s widely studied, leading to many high\sphinxhyphen{}quality multiparametric MRI datasets available for academic use. With these datasets, we are able to train some model on 3D volumetric data that will help us predict long\sphinxhyphen{}term patient outcomes.


\section{Existing Approaches}
\label{\detokenize{index:existing-approaches}}\begin{itemize}
\item {} 
\sphinxAtStartPar
\sphinxhref{https://www.nature.com/articles/s41598-025-97763-0}{Nature Scientific Reports (2025)}: Trained machine learning models to predict survival rates for patients undergoing breast cancer treatments. \sphinxstyleemphasis{(Note: Our project specifically builds on this domain by focusing on deep learning for 3D volumetric data and adding interpretability).}

\end{itemize}


\section{Datasets \& Methodology}
\label{\detokenize{index:datasets-methodology}}\begin{itemize}
\item {} 
\sphinxAtStartPar
\sphinxstylestrong{Dataset:} \sphinxhref{https://github.com/naomifridman/BreastDCEDL}{BreastDCEDL} A deep\sphinxhyphen{}learning\sphinxhyphen{}ready, labeled dataset that combines dynamic contrast\sphinxhyphen{}enhanced (DCE) MRI scans from three major clinical trials.

\item {} 
\sphinxAtStartPar
\sphinxstylestrong{Environment:} Google Colab.

\item {} 
\sphinxAtStartPar
\sphinxstylestrong{Architecture:} 3D Convolutional Neural Network (CNN) trained on volumetric MRI scans.

\item {} 
\sphinxAtStartPar
\sphinxstylestrong{Interpretability:} We will implement \sphinxstylestrong{HiResCAM} to visualize the model’s decision\sphinxhyphen{}making process. Because medical decisions require high transparency, HiResCAM will ensure the generated heatmaps are mathematically faithful to the model’s internal weights.

\item {} 
\sphinxAtStartPar
\sphinxstylestrong{Stretch Goal:} If time allows, we will implement custom positive/negative filters on the HiResCAM outputs. This toggle will allow clinicians to isolate the specific tissue features acting as positive contributors (evidence \sphinxstyleemphasis{for} pCR) versus negative contributors (evidence \sphinxstyleemphasis{against} pCR).

\end{itemize}


\section{Works Cited}
\label{\detokenize{index:works-cited}}
\sphinxAtStartPar
Draelos, Rachel Lea, and Lawrence Carin. “Use HiResCAM Instead of Grad\sphinxhyphen{}CAM for Faithful Explanations of Convolutional Neural Networks.” \sphinxstyleemphasis{arXiv}, 17 Nov. 2020, arxiv.org/abs/2011.08891.

\sphinxAtStartPar
Fridman, Naomi, et al. “BreastDCEDL: A Deep Learning\textendash{}Ready Breast DCE\sphinxhyphen{}MRI Dataset.” \sphinxstyleemphasis{Zenodo}, 9 June 2025, doi.org/10.5281/zenodo.15627233.

\sphinxAtStartPar
Toss, Angela, et al. “Predictive Factors for Relapse in Triple\sphinxhyphen{}Negative Breast Cancer Patients without Pathological Complete Response after Neoadjuvant Chemotherapy.” \sphinxstyleemphasis{Frontiers in Oncology}, vol. 12, 1 Dec. 2022, doi.org/10.3389/fonc.2022.1016295.

\sphinxstepscope


\chapter{Modules}
\label{\detokenize{modules:modules}}\label{\detokenize{modules::doc}}
\sphinxstepscope


\section{train}
\label{\detokenize{train:module-train}}\label{\detokenize{train:train}}\label{\detokenize{train::doc}}\index{module@\spxentry{module}!train@\spxentry{train}}\index{train@\spxentry{train}!module@\spxentry{module}}
\sphinxAtStartPar
Train a CNN using dataset.py
\index{ConvBlock (class in train)@\spxentry{ConvBlock}\spxextra{class in train}}

\begin{fulllineitems}
\phantomsection\label{\detokenize{train:train.ConvBlock}}
\pysigstartsignatures
\pysiglinewithargsret
{\sphinxbfcode{\sphinxupquote{\DUrole{k}{class}\DUrole{w}{ }}}\sphinxcode{\sphinxupquote{train.}}\sphinxbfcode{\sphinxupquote{ConvBlock}}}
{\sphinxparam{\DUrole{n}{num\_input\_channels}}\sphinxparamcomma \sphinxparam{\DUrole{n}{num\_output\_channels}}}
{}
\pysigstopsignatures
\sphinxAtStartPar
Bases: \sphinxcode{\sphinxupquote{Module}}
\index{forward() (train.ConvBlock method)@\spxentry{forward()}\spxextra{train.ConvBlock method}}

\begin{fulllineitems}
\phantomsection\label{\detokenize{train:train.ConvBlock.forward}}
\pysigstartsignatures
\pysiglinewithargsret
{\sphinxbfcode{\sphinxupquote{forward}}}
{\sphinxparam{\DUrole{n}{x}}}
{}
\pysigstopsignatures
\sphinxAtStartPar
Define the computation performed at every call.

\sphinxAtStartPar
Should be overridden by all subclasses.

\begin{sphinxadmonition}{note}{Note:}
\sphinxAtStartPar
Although the recipe for forward pass needs to be defined within
this function, one should call the \sphinxcode{\sphinxupquote{Module}} instance afterwards
instead of this since the former takes care of running the
registered hooks while the latter silently ignores them.
\end{sphinxadmonition}

\end{fulllineitems}


\end{fulllineitems}

\index{PcrCNN (class in train)@\spxentry{PcrCNN}\spxextra{class in train}}

\begin{fulllineitems}
\phantomsection\label{\detokenize{train:train.PcrCNN}}
\pysigstartsignatures
\pysiglinewithargsret
{\sphinxbfcode{\sphinxupquote{\DUrole{k}{class}\DUrole{w}{ }}}\sphinxcode{\sphinxupquote{train.}}\sphinxbfcode{\sphinxupquote{PcrCNN}}}
{\sphinxparam{\DUrole{n}{learning\_rate}\DUrole{o}{=}\DUrole{default_value}{0.0001}}\sphinxparamcomma \sphinxparam{\DUrole{n}{pos\_weight}\DUrole{o}{=}\DUrole{default_value}{1.0}}}
{}
\pysigstopsignatures
\sphinxAtStartPar
Bases: \sphinxcode{\sphinxupquote{LightningModule}}, \sphinxcode{\sphinxupquote{Module}}
\index{configure\_optimizers() (train.PcrCNN method)@\spxentry{configure\_optimizers()}\spxextra{train.PcrCNN method}}

\begin{fulllineitems}
\phantomsection\label{\detokenize{train:train.PcrCNN.configure_optimizers}}
\pysigstartsignatures
\pysiglinewithargsret
{\sphinxbfcode{\sphinxupquote{configure\_optimizers}}}
{}
{}
\pysigstopsignatures
\sphinxAtStartPar
Choose what optimizers and learning\sphinxhyphen{}rate schedulers to use in your optimization. Normally you’d need one.
But in the case of GANs or similar you might have multiple. Optimization with multiple optimizers only works in
the manual optimization mode.
\begin{quote}\begin{description}
\sphinxlineitem{Returns}
\sphinxAtStartPar

\sphinxAtStartPar
Any of these 6 options.
\begin{itemize}
\item {} 
\sphinxAtStartPar
\sphinxstylestrong{Single optimizer}.

\item {} 
\sphinxAtStartPar
\sphinxstylestrong{List or Tuple} of optimizers.

\item {} 
\sphinxAtStartPar
\sphinxstylestrong{Two lists} \sphinxhyphen{} The first list has multiple optimizers, and the second has multiple LR schedulers
(or multiple \sphinxcode{\sphinxupquote{lr\_scheduler\_config}}).

\item {} 
\sphinxAtStartPar
\sphinxstylestrong{Dictionary}, with an \sphinxcode{\sphinxupquote{"optimizer"}} key, and (optionally) a \sphinxcode{\sphinxupquote{"lr\_scheduler"}}
key whose value is a single LR scheduler or \sphinxcode{\sphinxupquote{lr\_scheduler\_config}}.

\item {} 
\sphinxAtStartPar
\sphinxstylestrong{None} \sphinxhyphen{} Fit will run without any optimizer.

\end{itemize}


\end{description}\end{quote}

\sphinxAtStartPar
The \sphinxcode{\sphinxupquote{lr\_scheduler\_config}} is a dictionary which contains the scheduler and its associated configuration.
The default configuration is shown below.

\begin{sphinxVerbatim}[commandchars=\\\{\}]
\PYG{n}{lr\PYGZus{}scheduler\PYGZus{}config} \PYG{o}{=} \PYG{p}{\PYGZob{}}
    \PYG{c+c1}{\PYGZsh{} REQUIRED: The scheduler instance}
    \PYG{l+s+s2}{\PYGZdq{}}\PYG{l+s+s2}{scheduler}\PYG{l+s+s2}{\PYGZdq{}}\PYG{p}{:} \PYG{n}{lr\PYGZus{}scheduler}\PYG{p}{,}
    \PYG{c+c1}{\PYGZsh{} The unit of the scheduler\PYGZsq{}s step size, could also be \PYGZsq{}step\PYGZsq{}.}
    \PYG{c+c1}{\PYGZsh{} \PYGZsq{}epoch\PYGZsq{} updates the scheduler on epoch end whereas \PYGZsq{}step\PYGZsq{}}
    \PYG{c+c1}{\PYGZsh{} updates it after a optimizer update.}
    \PYG{l+s+s2}{\PYGZdq{}}\PYG{l+s+s2}{interval}\PYG{l+s+s2}{\PYGZdq{}}\PYG{p}{:} \PYG{l+s+s2}{\PYGZdq{}}\PYG{l+s+s2}{epoch}\PYG{l+s+s2}{\PYGZdq{}}\PYG{p}{,}
    \PYG{c+c1}{\PYGZsh{} How many epochs/steps should pass between calls to}
    \PYG{c+c1}{\PYGZsh{} `scheduler.step()`. 1 corresponds to updating the learning}
    \PYG{c+c1}{\PYGZsh{} rate after every epoch/step.}
    \PYG{l+s+s2}{\PYGZdq{}}\PYG{l+s+s2}{frequency}\PYG{l+s+s2}{\PYGZdq{}}\PYG{p}{:} \PYG{l+m+mi}{1}\PYG{p}{,}
    \PYG{c+c1}{\PYGZsh{} Metric to monitor for schedulers like `ReduceLROnPlateau`}
    \PYG{l+s+s2}{\PYGZdq{}}\PYG{l+s+s2}{monitor}\PYG{l+s+s2}{\PYGZdq{}}\PYG{p}{:} \PYG{l+s+s2}{\PYGZdq{}}\PYG{l+s+s2}{val\PYGZus{}loss}\PYG{l+s+s2}{\PYGZdq{}}\PYG{p}{,}
    \PYG{c+c1}{\PYGZsh{} If set to `True`, will enforce that the value specified \PYGZsq{}monitor\PYGZsq{}}
    \PYG{c+c1}{\PYGZsh{} is available when the scheduler is updated, thus stopping}
    \PYG{c+c1}{\PYGZsh{} training if not found. If set to `False`, it will only produce a warning}
    \PYG{l+s+s2}{\PYGZdq{}}\PYG{l+s+s2}{strict}\PYG{l+s+s2}{\PYGZdq{}}\PYG{p}{:} \PYG{k+kc}{True}\PYG{p}{,}
    \PYG{c+c1}{\PYGZsh{} If using the `LearningRateMonitor` callback to monitor the}
    \PYG{c+c1}{\PYGZsh{} learning rate progress, this keyword can be used to specify}
    \PYG{c+c1}{\PYGZsh{} a custom logged name}
    \PYG{l+s+s2}{\PYGZdq{}}\PYG{l+s+s2}{name}\PYG{l+s+s2}{\PYGZdq{}}\PYG{p}{:} \PYG{k+kc}{None}\PYG{p}{,}
\PYG{p}{\PYGZcb{}}
\end{sphinxVerbatim}

\sphinxAtStartPar
When there are schedulers in which the \sphinxcode{\sphinxupquote{.step()}} method is conditioned on a value, such as the
\sphinxcode{\sphinxupquote{torch.optim.lr\_scheduler.ReduceLROnPlateau}} scheduler, Lightning requires that the
\sphinxcode{\sphinxupquote{lr\_scheduler\_config}} contains the keyword \sphinxcode{\sphinxupquote{"monitor"}} set to the metric name that the scheduler
should be conditioned on.

\sphinxAtStartPar
Metrics can be made available to monitor by simply logging it using
\sphinxcode{\sphinxupquote{self.log(\textquotesingle{}metric\_to\_track\textquotesingle{}, metric\_val)}} in your \sphinxcode{\sphinxupquote{LightningModule}}.

\begin{sphinxadmonition}{note}{Note:}
\sphinxAtStartPar
Some things to know:
\begin{itemize}
\item {} 
\sphinxAtStartPar
Lightning calls \sphinxcode{\sphinxupquote{.backward()}} and \sphinxcode{\sphinxupquote{.step()}} automatically in case of automatic optimization.

\item {} 
\sphinxAtStartPar
If a learning rate scheduler is specified in \sphinxcode{\sphinxupquote{configure\_optimizers()}} with key
\sphinxcode{\sphinxupquote{"interval"}} (default “epoch”) in the scheduler configuration, Lightning will call
the scheduler’s \sphinxcode{\sphinxupquote{.step()}} method automatically in case of automatic optimization.

\item {} 
\sphinxAtStartPar
If you use 16\sphinxhyphen{}bit precision (\sphinxcode{\sphinxupquote{precision=16}}), Lightning will automatically handle the optimizer.

\item {} 
\sphinxAtStartPar
If you use \sphinxcode{\sphinxupquote{torch.optim.LBFGS}}, Lightning handles the closure function automatically for you.

\item {} 
\sphinxAtStartPar
If you use multiple optimizers, you will have to switch to ‘manual optimization’ mode and step them
yourself.

\item {} 
\sphinxAtStartPar
If you need to control how often the optimizer steps, override the \sphinxcode{\sphinxupquote{optimizer\_step()}} hook.

\end{itemize}
\end{sphinxadmonition}

\end{fulllineitems}

\index{forward() (train.PcrCNN method)@\spxentry{forward()}\spxextra{train.PcrCNN method}}

\begin{fulllineitems}
\phantomsection\label{\detokenize{train:train.PcrCNN.forward}}
\pysigstartsignatures
\pysiglinewithargsret
{\sphinxbfcode{\sphinxupquote{forward}}}
{\sphinxparam{\DUrole{n}{x}}}
{}
\pysigstopsignatures
\sphinxAtStartPar
Same as \sphinxcode{\sphinxupquote{torch.nn.Module.forward()}}.
\begin{quote}\begin{description}
\sphinxlineitem{Parameters}\begin{itemize}
\item {} 
\sphinxAtStartPar
\sphinxstyleliteralstrong{\sphinxupquote{*args}} \textendash{} Whatever you decide to pass into the forward method.

\item {} 
\sphinxAtStartPar
\sphinxstyleliteralstrong{\sphinxupquote{**kwargs}} \textendash{} Keyword arguments are also possible.

\end{itemize}

\sphinxlineitem{Returns}
\sphinxAtStartPar
Your model’s output

\end{description}\end{quote}

\end{fulllineitems}

\index{on\_test\_epoch\_end() (train.PcrCNN method)@\spxentry{on\_test\_epoch\_end()}\spxextra{train.PcrCNN method}}

\begin{fulllineitems}
\phantomsection\label{\detokenize{train:train.PcrCNN.on_test_epoch_end}}
\pysigstartsignatures
\pysiglinewithargsret
{\sphinxbfcode{\sphinxupquote{on\_test\_epoch\_end}}}
{}
{}
\pysigstopsignatures
\sphinxAtStartPar
Called in the test loop at the very end of the epoch.

\end{fulllineitems}

\index{on\_train\_batch\_end() (train.PcrCNN method)@\spxentry{on\_train\_batch\_end()}\spxextra{train.PcrCNN method}}

\begin{fulllineitems}
\phantomsection\label{\detokenize{train:train.PcrCNN.on_train_batch_end}}
\pysigstartsignatures
\pysiglinewithargsret
{\sphinxbfcode{\sphinxupquote{on\_train\_batch\_end}}}
{\sphinxparam{\DUrole{n}{outputs}}\sphinxparamcomma \sphinxparam{\DUrole{n}{batch}}\sphinxparamcomma \sphinxparam{\DUrole{n}{batch\_idx}}}
{}
\pysigstopsignatures
\sphinxAtStartPar
Checks for NaN weights at the end of every training batch.

\end{fulllineitems}

\index{on\_train\_epoch\_end() (train.PcrCNN method)@\spxentry{on\_train\_epoch\_end()}\spxextra{train.PcrCNN method}}

\begin{fulllineitems}
\phantomsection\label{\detokenize{train:train.PcrCNN.on_train_epoch_end}}
\pysigstartsignatures
\pysiglinewithargsret
{\sphinxbfcode{\sphinxupquote{on\_train\_epoch\_end}}}
{}
{}
\pysigstopsignatures
\sphinxAtStartPar
Called in the training loop at the very end of the epoch.

\sphinxAtStartPar
To access all batch outputs at the end of the epoch, you can cache step outputs as an attribute of the
\sphinxcode{\sphinxupquote{LightningModule}} and access them in this hook:

\begin{sphinxVerbatim}[commandchars=\\\{\}]
\PYG{k}{class}\PYG{+w}{ }\PYG{n+nc}{MyLightningModule}\PYG{p}{(}\PYG{n}{L}\PYG{o}{.}\PYG{n}{LightningModule}\PYG{p}{)}\PYG{p}{:}
    \PYG{k}{def}\PYG{+w}{ }\PYG{n+nf+fm}{\PYGZus{}\PYGZus{}init\PYGZus{}\PYGZus{}}\PYG{p}{(}\PYG{n+nb+bp}{self}\PYG{p}{)}\PYG{p}{:}
        \PYG{n+nb}{super}\PYG{p}{(}\PYG{p}{)}\PYG{o}{.}\PYG{n+nf+fm}{\PYGZus{}\PYGZus{}init\PYGZus{}\PYGZus{}}\PYG{p}{(}\PYG{p}{)}
        \PYG{n+nb+bp}{self}\PYG{o}{.}\PYG{n}{training\PYGZus{}step\PYGZus{}outputs} \PYG{o}{=} \PYG{p}{[}\PYG{p}{]}

    \PYG{k}{def}\PYG{+w}{ }\PYG{n+nf}{training\PYGZus{}step}\PYG{p}{(}\PYG{n+nb+bp}{self}\PYG{p}{)}\PYG{p}{:}
        \PYG{n}{loss} \PYG{o}{=} \PYG{o}{.}\PYG{o}{.}\PYG{o}{.}
        \PYG{n+nb+bp}{self}\PYG{o}{.}\PYG{n}{training\PYGZus{}step\PYGZus{}outputs}\PYG{o}{.}\PYG{n}{append}\PYG{p}{(}\PYG{n}{loss}\PYG{p}{)}
        \PYG{k}{return} \PYG{n}{loss}

    \PYG{k}{def}\PYG{+w}{ }\PYG{n+nf}{on\PYGZus{}train\PYGZus{}epoch\PYGZus{}end}\PYG{p}{(}\PYG{n+nb+bp}{self}\PYG{p}{)}\PYG{p}{:}
        \PYG{c+c1}{\PYGZsh{} do something with all training\PYGZus{}step outputs, for example:}
        \PYG{n}{epoch\PYGZus{}mean} \PYG{o}{=} \PYG{n}{torch}\PYG{o}{.}\PYG{n}{stack}\PYG{p}{(}\PYG{n+nb+bp}{self}\PYG{o}{.}\PYG{n}{training\PYGZus{}step\PYGZus{}outputs}\PYG{p}{)}\PYG{o}{.}\PYG{n}{mean}\PYG{p}{(}\PYG{p}{)}
        \PYG{n+nb+bp}{self}\PYG{o}{.}\PYG{n}{log}\PYG{p}{(}\PYG{l+s+s2}{\PYGZdq{}}\PYG{l+s+s2}{training\PYGZus{}epoch\PYGZus{}mean}\PYG{l+s+s2}{\PYGZdq{}}\PYG{p}{,} \PYG{n}{epoch\PYGZus{}mean}\PYG{p}{)}
        \PYG{c+c1}{\PYGZsh{} free up the memory}
        \PYG{n+nb+bp}{self}\PYG{o}{.}\PYG{n}{training\PYGZus{}step\PYGZus{}outputs}\PYG{o}{.}\PYG{n}{clear}\PYG{p}{(}\PYG{p}{)}
\end{sphinxVerbatim}

\end{fulllineitems}

\index{on\_validation\_epoch\_end() (train.PcrCNN method)@\spxentry{on\_validation\_epoch\_end()}\spxextra{train.PcrCNN method}}

\begin{fulllineitems}
\phantomsection\label{\detokenize{train:train.PcrCNN.on_validation_epoch_end}}
\pysigstartsignatures
\pysiglinewithargsret
{\sphinxbfcode{\sphinxupquote{on\_validation\_epoch\_end}}}
{}
{}
\pysigstopsignatures
\sphinxAtStartPar
Called in the validation loop at the very end of the epoch.

\end{fulllineitems}

\index{shared\_step() (train.PcrCNN method)@\spxentry{shared\_step()}\spxextra{train.PcrCNN method}}

\begin{fulllineitems}
\phantomsection\label{\detokenize{train:train.PcrCNN.shared_step}}
\pysigstartsignatures
\pysiglinewithargsret
{\sphinxbfcode{\sphinxupquote{shared\_step}}}
{\sphinxparam{\DUrole{n}{batch}}}
{}
\pysigstopsignatures
\end{fulllineitems}

\index{test\_step() (train.PcrCNN method)@\spxentry{test\_step()}\spxextra{train.PcrCNN method}}

\begin{fulllineitems}
\phantomsection\label{\detokenize{train:train.PcrCNN.test_step}}
\pysigstartsignatures
\pysiglinewithargsret
{\sphinxbfcode{\sphinxupquote{test\_step}}}
{\sphinxparam{\DUrole{n}{batch}}\sphinxparamcomma \sphinxparam{\DUrole{n}{batch\_idx}}}
{}
\pysigstopsignatures
\sphinxAtStartPar
Operates on a single batch of data from the test set. In this step you’d normally generate examples or
calculate anything of interest such as accuracy.
\begin{quote}\begin{description}
\sphinxlineitem{Parameters}\begin{itemize}
\item {} 
\sphinxAtStartPar
\sphinxstyleliteralstrong{\sphinxupquote{batch}} \textendash{} The output of your data iterable, normally a \sphinxcode{\sphinxupquote{DataLoader}}.

\item {} 
\sphinxAtStartPar
\sphinxstyleliteralstrong{\sphinxupquote{batch\_idx}} \textendash{} The index of this batch.

\item {} 
\sphinxAtStartPar
\sphinxstyleliteralstrong{\sphinxupquote{dataloader\_idx}} \textendash{} The index of the dataloader that produced this batch.
(only if multiple dataloaders used)

\end{itemize}

\sphinxlineitem{Returns}
\sphinxAtStartPar
\begin{itemize}
\item {} 
\sphinxAtStartPar
\sphinxcode{\sphinxupquote{Tensor}} \sphinxhyphen{} The loss tensor

\item {} 
\sphinxAtStartPar
\sphinxcode{\sphinxupquote{dict}} \sphinxhyphen{} A dictionary. Can include any keys, but must include the key \sphinxcode{\sphinxupquote{\textquotesingle{}loss\textquotesingle{}}}.

\item {} 
\sphinxAtStartPar
\sphinxcode{\sphinxupquote{None}} \sphinxhyphen{} Skip to the next batch.

\end{itemize}


\end{description}\end{quote}

\begin{sphinxVerbatim}[commandchars=\\\{\}]
\PYG{c+c1}{\PYGZsh{} if you have one test dataloader:}
\PYG{k}{def}\PYG{+w}{ }\PYG{n+nf}{test\PYGZus{}step}\PYG{p}{(}\PYG{n+nb+bp}{self}\PYG{p}{,} \PYG{n}{batch}\PYG{p}{,} \PYG{n}{batch\PYGZus{}idx}\PYG{p}{)}\PYG{p}{:} \PYG{o}{.}\PYG{o}{.}\PYG{o}{.}


\PYG{c+c1}{\PYGZsh{} if you have multiple test dataloaders:}
\PYG{k}{def}\PYG{+w}{ }\PYG{n+nf}{test\PYGZus{}step}\PYG{p}{(}\PYG{n+nb+bp}{self}\PYG{p}{,} \PYG{n}{batch}\PYG{p}{,} \PYG{n}{batch\PYGZus{}idx}\PYG{p}{,} \PYG{n}{dataloader\PYGZus{}idx}\PYG{o}{=}\PYG{l+m+mi}{0}\PYG{p}{)}\PYG{p}{:} \PYG{o}{.}\PYG{o}{.}\PYG{o}{.}
\end{sphinxVerbatim}

\sphinxAtStartPar
Examples:

\begin{sphinxVerbatim}[commandchars=\\\{\}]
\PYG{c+c1}{\PYGZsh{} CASE 1: A single test dataset}
\PYG{k}{def}\PYG{+w}{ }\PYG{n+nf}{test\PYGZus{}step}\PYG{p}{(}\PYG{n+nb+bp}{self}\PYG{p}{,} \PYG{n}{batch}\PYG{p}{,} \PYG{n}{batch\PYGZus{}idx}\PYG{p}{)}\PYG{p}{:}
    \PYG{n}{x}\PYG{p}{,} \PYG{n}{y} \PYG{o}{=} \PYG{n}{batch}

    \PYG{c+c1}{\PYGZsh{} implement your own}
    \PYG{n}{out} \PYG{o}{=} \PYG{n+nb+bp}{self}\PYG{p}{(}\PYG{n}{x}\PYG{p}{)}
    \PYG{n}{loss} \PYG{o}{=} \PYG{n+nb+bp}{self}\PYG{o}{.}\PYG{n}{loss}\PYG{p}{(}\PYG{n}{out}\PYG{p}{,} \PYG{n}{y}\PYG{p}{)}

    \PYG{c+c1}{\PYGZsh{} log 6 example images}
    \PYG{c+c1}{\PYGZsh{} or generated text... or whatever}
    \PYG{n}{sample\PYGZus{}imgs} \PYG{o}{=} \PYG{n}{x}\PYG{p}{[}\PYG{p}{:}\PYG{l+m+mi}{6}\PYG{p}{]}
    \PYG{n}{grid} \PYG{o}{=} \PYG{n}{torchvision}\PYG{o}{.}\PYG{n}{utils}\PYG{o}{.}\PYG{n}{make\PYGZus{}grid}\PYG{p}{(}\PYG{n}{sample\PYGZus{}imgs}\PYG{p}{)}
    \PYG{n+nb+bp}{self}\PYG{o}{.}\PYG{n}{logger}\PYG{o}{.}\PYG{n}{experiment}\PYG{o}{.}\PYG{n}{add\PYGZus{}image}\PYG{p}{(}\PYG{l+s+s1}{\PYGZsq{}}\PYG{l+s+s1}{example\PYGZus{}images}\PYG{l+s+s1}{\PYGZsq{}}\PYG{p}{,} \PYG{n}{grid}\PYG{p}{,} \PYG{l+m+mi}{0}\PYG{p}{)}

    \PYG{c+c1}{\PYGZsh{} calculate acc}
    \PYG{n}{labels\PYGZus{}hat} \PYG{o}{=} \PYG{n}{torch}\PYG{o}{.}\PYG{n}{argmax}\PYG{p}{(}\PYG{n}{out}\PYG{p}{,} \PYG{n}{dim}\PYG{o}{=}\PYG{l+m+mi}{1}\PYG{p}{)}
    \PYG{n}{test\PYGZus{}acc} \PYG{o}{=} \PYG{n}{torch}\PYG{o}{.}\PYG{n}{sum}\PYG{p}{(}\PYG{n}{y} \PYG{o}{==} \PYG{n}{labels\PYGZus{}hat}\PYG{p}{)}\PYG{o}{.}\PYG{n}{item}\PYG{p}{(}\PYG{p}{)} \PYG{o}{/} \PYG{p}{(}\PYG{n+nb}{len}\PYG{p}{(}\PYG{n}{y}\PYG{p}{)} \PYG{o}{*} \PYG{l+m+mf}{1.0}\PYG{p}{)}

    \PYG{c+c1}{\PYGZsh{} log the outputs!}
    \PYG{n+nb+bp}{self}\PYG{o}{.}\PYG{n}{log\PYGZus{}dict}\PYG{p}{(}\PYG{p}{\PYGZob{}}\PYG{l+s+s1}{\PYGZsq{}}\PYG{l+s+s1}{test\PYGZus{}loss}\PYG{l+s+s1}{\PYGZsq{}}\PYG{p}{:} \PYG{n}{loss}\PYG{p}{,} \PYG{l+s+s1}{\PYGZsq{}}\PYG{l+s+s1}{test\PYGZus{}acc}\PYG{l+s+s1}{\PYGZsq{}}\PYG{p}{:} \PYG{n}{test\PYGZus{}acc}\PYG{p}{\PYGZcb{}}\PYG{p}{)}
\end{sphinxVerbatim}

\sphinxAtStartPar
If you pass in multiple test dataloaders, {\hyperref[\detokenize{train:train.PcrCNN.test_step}]{\sphinxcrossref{\sphinxcode{\sphinxupquote{test\_step()}}}}} will have an additional argument. We recommend
setting the default value of 0 so that you can quickly switch between single and multiple dataloaders.

\begin{sphinxVerbatim}[commandchars=\\\{\}]
\PYG{c+c1}{\PYGZsh{} CASE 2: multiple test dataloaders}
\PYG{k}{def}\PYG{+w}{ }\PYG{n+nf}{test\PYGZus{}step}\PYG{p}{(}\PYG{n+nb+bp}{self}\PYG{p}{,} \PYG{n}{batch}\PYG{p}{,} \PYG{n}{batch\PYGZus{}idx}\PYG{p}{,} \PYG{n}{dataloader\PYGZus{}idx}\PYG{o}{=}\PYG{l+m+mi}{0}\PYG{p}{)}\PYG{p}{:}
    \PYG{c+c1}{\PYGZsh{} dataloader\PYGZus{}idx tells you which dataset this is.}
    \PYG{n}{x}\PYG{p}{,} \PYG{n}{y} \PYG{o}{=} \PYG{n}{batch}

    \PYG{c+c1}{\PYGZsh{} implement your own}
    \PYG{n}{out} \PYG{o}{=} \PYG{n+nb+bp}{self}\PYG{p}{(}\PYG{n}{x}\PYG{p}{)}

    \PYG{k}{if} \PYG{n}{dataloader\PYGZus{}idx} \PYG{o}{==} \PYG{l+m+mi}{0}\PYG{p}{:}
        \PYG{n}{loss} \PYG{o}{=} \PYG{n+nb+bp}{self}\PYG{o}{.}\PYG{n}{loss0}\PYG{p}{(}\PYG{n}{out}\PYG{p}{,} \PYG{n}{y}\PYG{p}{)}
    \PYG{k}{else}\PYG{p}{:}
        \PYG{n}{loss} \PYG{o}{=} \PYG{n+nb+bp}{self}\PYG{o}{.}\PYG{n}{loss1}\PYG{p}{(}\PYG{n}{out}\PYG{p}{,} \PYG{n}{y}\PYG{p}{)}

    \PYG{c+c1}{\PYGZsh{} calculate acc}
    \PYG{n}{labels\PYGZus{}hat} \PYG{o}{=} \PYG{n}{torch}\PYG{o}{.}\PYG{n}{argmax}\PYG{p}{(}\PYG{n}{out}\PYG{p}{,} \PYG{n}{dim}\PYG{o}{=}\PYG{l+m+mi}{1}\PYG{p}{)}
    \PYG{n}{acc} \PYG{o}{=} \PYG{n}{torch}\PYG{o}{.}\PYG{n}{sum}\PYG{p}{(}\PYG{n}{y} \PYG{o}{==} \PYG{n}{labels\PYGZus{}hat}\PYG{p}{)}\PYG{o}{.}\PYG{n}{item}\PYG{p}{(}\PYG{p}{)} \PYG{o}{/} \PYG{p}{(}\PYG{n+nb}{len}\PYG{p}{(}\PYG{n}{y}\PYG{p}{)} \PYG{o}{*} \PYG{l+m+mf}{1.0}\PYG{p}{)}

    \PYG{c+c1}{\PYGZsh{} log the outputs separately for each dataloader}
    \PYG{n+nb+bp}{self}\PYG{o}{.}\PYG{n}{log\PYGZus{}dict}\PYG{p}{(}\PYG{p}{\PYGZob{}}\PYG{l+s+sa}{f}\PYG{l+s+s2}{\PYGZdq{}}\PYG{l+s+s2}{test\PYGZus{}loss\PYGZus{}}\PYG{l+s+si}{\PYGZob{}}\PYG{n}{dataloader\PYGZus{}idx}\PYG{l+s+si}{\PYGZcb{}}\PYG{l+s+s2}{\PYGZdq{}}\PYG{p}{:} \PYG{n}{loss}\PYG{p}{,} \PYG{l+s+sa}{f}\PYG{l+s+s2}{\PYGZdq{}}\PYG{l+s+s2}{test\PYGZus{}acc\PYGZus{}}\PYG{l+s+si}{\PYGZob{}}\PYG{n}{dataloader\PYGZus{}idx}\PYG{l+s+si}{\PYGZcb{}}\PYG{l+s+s2}{\PYGZdq{}}\PYG{p}{:} \PYG{n}{acc}\PYG{p}{\PYGZcb{}}\PYG{p}{)}
\end{sphinxVerbatim}

\begin{sphinxadmonition}{note}{Note:}
\sphinxAtStartPar
If you don’t need to test you don’t need to implement this method.
\end{sphinxadmonition}

\begin{sphinxadmonition}{note}{Note:}
\sphinxAtStartPar
When the {\hyperref[\detokenize{train:train.PcrCNN.test_step}]{\sphinxcrossref{\sphinxcode{\sphinxupquote{test\_step()}}}}} is called, the model has been put in eval mode and
PyTorch gradients have been disabled. At the end of the test epoch, the model goes back
to training mode and gradients are enabled.
\end{sphinxadmonition}

\end{fulllineitems}

\index{training\_step() (train.PcrCNN method)@\spxentry{training\_step()}\spxextra{train.PcrCNN method}}

\begin{fulllineitems}
\phantomsection\label{\detokenize{train:train.PcrCNN.training_step}}
\pysigstartsignatures
\pysiglinewithargsret
{\sphinxbfcode{\sphinxupquote{training\_step}}}
{\sphinxparam{\DUrole{n}{batch}}\sphinxparamcomma \sphinxparam{\DUrole{n}{batch\_idx}}}
{}
\pysigstopsignatures
\sphinxAtStartPar
Here you compute and return the training loss and some additional metrics for e.g. the progress bar or
logger.
\begin{quote}\begin{description}
\sphinxlineitem{Parameters}\begin{itemize}
\item {} 
\sphinxAtStartPar
\sphinxstyleliteralstrong{\sphinxupquote{batch}} \textendash{} The output of your data iterable, normally a \sphinxcode{\sphinxupquote{DataLoader}}.

\item {} 
\sphinxAtStartPar
\sphinxstyleliteralstrong{\sphinxupquote{batch\_idx}} \textendash{} The index of this batch.

\item {} 
\sphinxAtStartPar
\sphinxstyleliteralstrong{\sphinxupquote{dataloader\_idx}} \textendash{} The index of the dataloader that produced this batch.
(only if multiple dataloaders used)

\end{itemize}

\sphinxlineitem{Returns}
\sphinxAtStartPar
\begin{itemize}
\item {} 
\sphinxAtStartPar
\sphinxcode{\sphinxupquote{Tensor}} \sphinxhyphen{} The loss tensor

\item {} 
\sphinxAtStartPar
\sphinxcode{\sphinxupquote{dict}} \sphinxhyphen{} A dictionary which can include any keys, but must include the key \sphinxcode{\sphinxupquote{\textquotesingle{}loss\textquotesingle{}}} in the case of
automatic optimization.

\item {} 
\sphinxAtStartPar
\sphinxcode{\sphinxupquote{None}} \sphinxhyphen{} In automatic optimization, this will skip to the next batch (but is not supported for
multi\sphinxhyphen{}GPU, TPU, or DeepSpeed). For manual optimization, this has no special meaning, as returning
the loss is not required.

\end{itemize}


\end{description}\end{quote}

\sphinxAtStartPar
In this step you’d normally do the forward pass and calculate the loss for a batch.
You can also do fancier things like multiple forward passes or something model specific.

\sphinxAtStartPar
Example:

\begin{sphinxVerbatim}[commandchars=\\\{\}]
\PYG{k}{def}\PYG{+w}{ }\PYG{n+nf}{training\PYGZus{}step}\PYG{p}{(}\PYG{n+nb+bp}{self}\PYG{p}{,} \PYG{n}{batch}\PYG{p}{,} \PYG{n}{batch\PYGZus{}idx}\PYG{p}{)}\PYG{p}{:}
    \PYG{n}{x}\PYG{p}{,} \PYG{n}{y}\PYG{p}{,} \PYG{n}{z} \PYG{o}{=} \PYG{n}{batch}
    \PYG{n}{out} \PYG{o}{=} \PYG{n+nb+bp}{self}\PYG{o}{.}\PYG{n}{encoder}\PYG{p}{(}\PYG{n}{x}\PYG{p}{)}
    \PYG{n}{loss} \PYG{o}{=} \PYG{n+nb+bp}{self}\PYG{o}{.}\PYG{n}{loss}\PYG{p}{(}\PYG{n}{out}\PYG{p}{,} \PYG{n}{x}\PYG{p}{)}
    \PYG{k}{return} \PYG{n}{loss}
\end{sphinxVerbatim}

\sphinxAtStartPar
To use multiple optimizers, you can switch to ‘manual optimization’ and control their stepping:

\begin{sphinxVerbatim}[commandchars=\\\{\}]
\PYG{k}{def}\PYG{+w}{ }\PYG{n+nf+fm}{\PYGZus{}\PYGZus{}init\PYGZus{}\PYGZus{}}\PYG{p}{(}\PYG{n+nb+bp}{self}\PYG{p}{)}\PYG{p}{:}
    \PYG{n+nb}{super}\PYG{p}{(}\PYG{p}{)}\PYG{o}{.}\PYG{n+nf+fm}{\PYGZus{}\PYGZus{}init\PYGZus{}\PYGZus{}}\PYG{p}{(}\PYG{p}{)}
    \PYG{n+nb+bp}{self}\PYG{o}{.}\PYG{n}{automatic\PYGZus{}optimization} \PYG{o}{=} \PYG{k+kc}{False}


\PYG{c+c1}{\PYGZsh{} Multiple optimizers (e.g.: GANs)}
\PYG{k}{def}\PYG{+w}{ }\PYG{n+nf}{training\PYGZus{}step}\PYG{p}{(}\PYG{n+nb+bp}{self}\PYG{p}{,} \PYG{n}{batch}\PYG{p}{,} \PYG{n}{batch\PYGZus{}idx}\PYG{p}{)}\PYG{p}{:}
    \PYG{n}{opt1}\PYG{p}{,} \PYG{n}{opt2} \PYG{o}{=} \PYG{n+nb+bp}{self}\PYG{o}{.}\PYG{n}{optimizers}\PYG{p}{(}\PYG{p}{)}

    \PYG{c+c1}{\PYGZsh{} do training\PYGZus{}step with encoder}
    \PYG{o}{.}\PYG{o}{.}\PYG{o}{.}
    \PYG{n}{opt1}\PYG{o}{.}\PYG{n}{step}\PYG{p}{(}\PYG{p}{)}
    \PYG{c+c1}{\PYGZsh{} do training\PYGZus{}step with decoder}
    \PYG{o}{.}\PYG{o}{.}\PYG{o}{.}
    \PYG{n}{opt2}\PYG{o}{.}\PYG{n}{step}\PYG{p}{(}\PYG{p}{)}
\end{sphinxVerbatim}

\begin{sphinxadmonition}{note}{Note:}
\sphinxAtStartPar
When \sphinxcode{\sphinxupquote{accumulate\_grad\_batches}} \textgreater{} 1, the loss returned here will be automatically
normalized by \sphinxcode{\sphinxupquote{accumulate\_grad\_batches}} internally.
\end{sphinxadmonition}

\end{fulllineitems}

\index{validation\_step() (train.PcrCNN method)@\spxentry{validation\_step()}\spxextra{train.PcrCNN method}}

\begin{fulllineitems}
\phantomsection\label{\detokenize{train:train.PcrCNN.validation_step}}
\pysigstartsignatures
\pysiglinewithargsret
{\sphinxbfcode{\sphinxupquote{validation\_step}}}
{\sphinxparam{\DUrole{n}{batch}}\sphinxparamcomma \sphinxparam{\DUrole{n}{batch\_idx}}}
{}
\pysigstopsignatures
\sphinxAtStartPar
Operates on a single batch of data from the validation set. In this step you’d might generate examples or
calculate anything of interest like accuracy.
\begin{quote}\begin{description}
\sphinxlineitem{Parameters}\begin{itemize}
\item {} 
\sphinxAtStartPar
\sphinxstyleliteralstrong{\sphinxupquote{batch}} \textendash{} The output of your data iterable, normally a \sphinxcode{\sphinxupquote{DataLoader}}.

\item {} 
\sphinxAtStartPar
\sphinxstyleliteralstrong{\sphinxupquote{batch\_idx}} \textendash{} The index of this batch.

\item {} 
\sphinxAtStartPar
\sphinxstyleliteralstrong{\sphinxupquote{dataloader\_idx}} \textendash{} The index of the dataloader that produced this batch.
(only if multiple dataloaders used)

\end{itemize}

\sphinxlineitem{Returns}
\sphinxAtStartPar
\begin{itemize}
\item {} 
\sphinxAtStartPar
\sphinxcode{\sphinxupquote{Tensor}} \sphinxhyphen{} The loss tensor

\item {} 
\sphinxAtStartPar
\sphinxcode{\sphinxupquote{dict}} \sphinxhyphen{} A dictionary. Can include any keys, but must include the key \sphinxcode{\sphinxupquote{\textquotesingle{}loss\textquotesingle{}}}.

\item {} 
\sphinxAtStartPar
\sphinxcode{\sphinxupquote{None}} \sphinxhyphen{} Skip to the next batch.

\end{itemize}


\end{description}\end{quote}

\begin{sphinxVerbatim}[commandchars=\\\{\}]
\PYG{c+c1}{\PYGZsh{} if you have one val dataloader:}
\PYG{k}{def}\PYG{+w}{ }\PYG{n+nf}{validation\PYGZus{}step}\PYG{p}{(}\PYG{n+nb+bp}{self}\PYG{p}{,} \PYG{n}{batch}\PYG{p}{,} \PYG{n}{batch\PYGZus{}idx}\PYG{p}{)}\PYG{p}{:} \PYG{o}{.}\PYG{o}{.}\PYG{o}{.}


\PYG{c+c1}{\PYGZsh{} if you have multiple val dataloaders:}
\PYG{k}{def}\PYG{+w}{ }\PYG{n+nf}{validation\PYGZus{}step}\PYG{p}{(}\PYG{n+nb+bp}{self}\PYG{p}{,} \PYG{n}{batch}\PYG{p}{,} \PYG{n}{batch\PYGZus{}idx}\PYG{p}{,} \PYG{n}{dataloader\PYGZus{}idx}\PYG{o}{=}\PYG{l+m+mi}{0}\PYG{p}{)}\PYG{p}{:} \PYG{o}{.}\PYG{o}{.}\PYG{o}{.}
\end{sphinxVerbatim}

\sphinxAtStartPar
Examples:

\begin{sphinxVerbatim}[commandchars=\\\{\}]
\PYG{c+c1}{\PYGZsh{} CASE 1: A single validation dataset}
\PYG{k}{def}\PYG{+w}{ }\PYG{n+nf}{validation\PYGZus{}step}\PYG{p}{(}\PYG{n+nb+bp}{self}\PYG{p}{,} \PYG{n}{batch}\PYG{p}{,} \PYG{n}{batch\PYGZus{}idx}\PYG{p}{)}\PYG{p}{:}
    \PYG{n}{x}\PYG{p}{,} \PYG{n}{y} \PYG{o}{=} \PYG{n}{batch}

    \PYG{c+c1}{\PYGZsh{} implement your own}
    \PYG{n}{out} \PYG{o}{=} \PYG{n+nb+bp}{self}\PYG{p}{(}\PYG{n}{x}\PYG{p}{)}
    \PYG{n}{loss} \PYG{o}{=} \PYG{n+nb+bp}{self}\PYG{o}{.}\PYG{n}{loss}\PYG{p}{(}\PYG{n}{out}\PYG{p}{,} \PYG{n}{y}\PYG{p}{)}

    \PYG{c+c1}{\PYGZsh{} log 6 example images}
    \PYG{c+c1}{\PYGZsh{} or generated text... or whatever}
    \PYG{n}{sample\PYGZus{}imgs} \PYG{o}{=} \PYG{n}{x}\PYG{p}{[}\PYG{p}{:}\PYG{l+m+mi}{6}\PYG{p}{]}
    \PYG{n}{grid} \PYG{o}{=} \PYG{n}{torchvision}\PYG{o}{.}\PYG{n}{utils}\PYG{o}{.}\PYG{n}{make\PYGZus{}grid}\PYG{p}{(}\PYG{n}{sample\PYGZus{}imgs}\PYG{p}{)}
    \PYG{n+nb+bp}{self}\PYG{o}{.}\PYG{n}{logger}\PYG{o}{.}\PYG{n}{experiment}\PYG{o}{.}\PYG{n}{add\PYGZus{}image}\PYG{p}{(}\PYG{l+s+s1}{\PYGZsq{}}\PYG{l+s+s1}{example\PYGZus{}images}\PYG{l+s+s1}{\PYGZsq{}}\PYG{p}{,} \PYG{n}{grid}\PYG{p}{,} \PYG{l+m+mi}{0}\PYG{p}{)}

    \PYG{c+c1}{\PYGZsh{} calculate acc}
    \PYG{n}{labels\PYGZus{}hat} \PYG{o}{=} \PYG{n}{torch}\PYG{o}{.}\PYG{n}{argmax}\PYG{p}{(}\PYG{n}{out}\PYG{p}{,} \PYG{n}{dim}\PYG{o}{=}\PYG{l+m+mi}{1}\PYG{p}{)}
    \PYG{n}{val\PYGZus{}acc} \PYG{o}{=} \PYG{n}{torch}\PYG{o}{.}\PYG{n}{sum}\PYG{p}{(}\PYG{n}{y} \PYG{o}{==} \PYG{n}{labels\PYGZus{}hat}\PYG{p}{)}\PYG{o}{.}\PYG{n}{item}\PYG{p}{(}\PYG{p}{)} \PYG{o}{/} \PYG{p}{(}\PYG{n+nb}{len}\PYG{p}{(}\PYG{n}{y}\PYG{p}{)} \PYG{o}{*} \PYG{l+m+mf}{1.0}\PYG{p}{)}

    \PYG{c+c1}{\PYGZsh{} log the outputs!}
    \PYG{n+nb+bp}{self}\PYG{o}{.}\PYG{n}{log\PYGZus{}dict}\PYG{p}{(}\PYG{p}{\PYGZob{}}\PYG{l+s+s1}{\PYGZsq{}}\PYG{l+s+s1}{val\PYGZus{}loss}\PYG{l+s+s1}{\PYGZsq{}}\PYG{p}{:} \PYG{n}{loss}\PYG{p}{,} \PYG{l+s+s1}{\PYGZsq{}}\PYG{l+s+s1}{val\PYGZus{}acc}\PYG{l+s+s1}{\PYGZsq{}}\PYG{p}{:} \PYG{n}{val\PYGZus{}acc}\PYG{p}{\PYGZcb{}}\PYG{p}{)}
\end{sphinxVerbatim}

\sphinxAtStartPar
If you pass in multiple val dataloaders, {\hyperref[\detokenize{train:train.PcrCNN.validation_step}]{\sphinxcrossref{\sphinxcode{\sphinxupquote{validation\_step()}}}}} will have an additional argument. We recommend
setting the default value of 0 so that you can quickly switch between single and multiple dataloaders.

\begin{sphinxVerbatim}[commandchars=\\\{\}]
\PYG{c+c1}{\PYGZsh{} CASE 2: multiple validation dataloaders}
\PYG{k}{def}\PYG{+w}{ }\PYG{n+nf}{validation\PYGZus{}step}\PYG{p}{(}\PYG{n+nb+bp}{self}\PYG{p}{,} \PYG{n}{batch}\PYG{p}{,} \PYG{n}{batch\PYGZus{}idx}\PYG{p}{,} \PYG{n}{dataloader\PYGZus{}idx}\PYG{o}{=}\PYG{l+m+mi}{0}\PYG{p}{)}\PYG{p}{:}
    \PYG{c+c1}{\PYGZsh{} dataloader\PYGZus{}idx tells you which dataset this is.}
    \PYG{n}{x}\PYG{p}{,} \PYG{n}{y} \PYG{o}{=} \PYG{n}{batch}

    \PYG{c+c1}{\PYGZsh{} implement your own}
    \PYG{n}{out} \PYG{o}{=} \PYG{n+nb+bp}{self}\PYG{p}{(}\PYG{n}{x}\PYG{p}{)}

    \PYG{k}{if} \PYG{n}{dataloader\PYGZus{}idx} \PYG{o}{==} \PYG{l+m+mi}{0}\PYG{p}{:}
        \PYG{n}{loss} \PYG{o}{=} \PYG{n+nb+bp}{self}\PYG{o}{.}\PYG{n}{loss0}\PYG{p}{(}\PYG{n}{out}\PYG{p}{,} \PYG{n}{y}\PYG{p}{)}
    \PYG{k}{else}\PYG{p}{:}
        \PYG{n}{loss} \PYG{o}{=} \PYG{n+nb+bp}{self}\PYG{o}{.}\PYG{n}{loss1}\PYG{p}{(}\PYG{n}{out}\PYG{p}{,} \PYG{n}{y}\PYG{p}{)}

    \PYG{c+c1}{\PYGZsh{} calculate acc}
    \PYG{n}{labels\PYGZus{}hat} \PYG{o}{=} \PYG{n}{torch}\PYG{o}{.}\PYG{n}{argmax}\PYG{p}{(}\PYG{n}{out}\PYG{p}{,} \PYG{n}{dim}\PYG{o}{=}\PYG{l+m+mi}{1}\PYG{p}{)}
    \PYG{n}{acc} \PYG{o}{=} \PYG{n}{torch}\PYG{o}{.}\PYG{n}{sum}\PYG{p}{(}\PYG{n}{y} \PYG{o}{==} \PYG{n}{labels\PYGZus{}hat}\PYG{p}{)}\PYG{o}{.}\PYG{n}{item}\PYG{p}{(}\PYG{p}{)} \PYG{o}{/} \PYG{p}{(}\PYG{n+nb}{len}\PYG{p}{(}\PYG{n}{y}\PYG{p}{)} \PYG{o}{*} \PYG{l+m+mf}{1.0}\PYG{p}{)}

    \PYG{c+c1}{\PYGZsh{} log the outputs separately for each dataloader}
    \PYG{n+nb+bp}{self}\PYG{o}{.}\PYG{n}{log\PYGZus{}dict}\PYG{p}{(}\PYG{p}{\PYGZob{}}\PYG{l+s+sa}{f}\PYG{l+s+s2}{\PYGZdq{}}\PYG{l+s+s2}{val\PYGZus{}loss\PYGZus{}}\PYG{l+s+si}{\PYGZob{}}\PYG{n}{dataloader\PYGZus{}idx}\PYG{l+s+si}{\PYGZcb{}}\PYG{l+s+s2}{\PYGZdq{}}\PYG{p}{:} \PYG{n}{loss}\PYG{p}{,} \PYG{l+s+sa}{f}\PYG{l+s+s2}{\PYGZdq{}}\PYG{l+s+s2}{val\PYGZus{}acc\PYGZus{}}\PYG{l+s+si}{\PYGZob{}}\PYG{n}{dataloader\PYGZus{}idx}\PYG{l+s+si}{\PYGZcb{}}\PYG{l+s+s2}{\PYGZdq{}}\PYG{p}{:} \PYG{n}{acc}\PYG{p}{\PYGZcb{}}\PYG{p}{)}
\end{sphinxVerbatim}

\begin{sphinxadmonition}{note}{Note:}
\sphinxAtStartPar
If you don’t need to validate you don’t need to implement this method.
\end{sphinxadmonition}

\begin{sphinxadmonition}{note}{Note:}
\sphinxAtStartPar
When the {\hyperref[\detokenize{train:train.PcrCNN.validation_step}]{\sphinxcrossref{\sphinxcode{\sphinxupquote{validation\_step()}}}}} is called, the model has been put in eval mode
and PyTorch gradients have been disabled. At the end of validation,
the model goes back to training mode and gradients are enabled.
\end{sphinxadmonition}

\end{fulllineitems}


\end{fulllineitems}

\index{main() (in module train)@\spxentry{main()}\spxextra{in module train}}

\begin{fulllineitems}
\phantomsection\label{\detokenize{train:train.main}}
\pysigstartsignatures
\pysiglinewithargsret
{\sphinxcode{\sphinxupquote{train.}}\sphinxbfcode{\sphinxupquote{main}}}
{}
{}
\pysigstopsignatures
\end{fulllineitems}


\sphinxstepscope


\section{dataset}
\label{\detokenize{dataset:module-dataset}}\label{\detokenize{dataset:dataset}}\label{\detokenize{dataset::doc}}\index{module@\spxentry{module}!dataset@\spxentry{dataset}}\index{dataset@\spxentry{dataset}!module@\spxentry{module}}
\sphinxAtStartPar
Download the min crop dataset from \sphinxurl{https://zenodo.org/records/18114231}.

\sphinxAtStartPar
DO NOT COMMIT DATASET.
\index{BreastDCEDataset (class in dataset)@\spxentry{BreastDCEDataset}\spxextra{class in dataset}}

\begin{fulllineitems}
\phantomsection\label{\detokenize{dataset:dataset.BreastDCEDataset}}
\pysigstartsignatures
\pysiglinewithargsret
{\sphinxbfcode{\sphinxupquote{\DUrole{k}{class}\DUrole{w}{ }}}\sphinxcode{\sphinxupquote{dataset.}}\sphinxbfcode{\sphinxupquote{BreastDCEDataset}}}
{\sphinxparam{\DUrole{n}{csv\_dir}}\sphinxparamcomma \sphinxparam{\DUrole{n}{data\_dir}}\sphinxparamcomma \sphinxparam{\DUrole{n}{split}\DUrole{o}{=}\DUrole{default_value}{Split.TRAIN}}}
{}
\pysigstopsignatures
\sphinxAtStartPar
Bases: \sphinxcode{\sphinxupquote{Dataset}}

\end{fulllineitems}

\index{Split (class in dataset)@\spxentry{Split}\spxextra{class in dataset}}

\begin{fulllineitems}
\phantomsection\label{\detokenize{dataset:dataset.Split}}
\pysigstartsignatures
\pysiglinewithargsret
{\sphinxbfcode{\sphinxupquote{\DUrole{k}{class}\DUrole{w}{ }}}\sphinxcode{\sphinxupquote{dataset.}}\sphinxbfcode{\sphinxupquote{Split}}}
{\sphinxparam{\DUrole{o}{*}\DUrole{n}{values}}}
{}
\pysigstopsignatures
\sphinxAtStartPar
Bases: \sphinxcode{\sphinxupquote{IntEnum}}
\index{TEST (dataset.Split attribute)@\spxentry{TEST}\spxextra{dataset.Split attribute}}

\begin{fulllineitems}
\phantomsection\label{\detokenize{dataset:dataset.Split.TEST}}
\pysigstartsignatures
\pysigline
{\sphinxbfcode{\sphinxupquote{TEST}}\sphinxbfcode{\sphinxupquote{\DUrole{w}{ }\DUrole{p}{=}\DUrole{w}{ }2}}}
\pysigstopsignatures
\end{fulllineitems}

\index{TRAIN (dataset.Split attribute)@\spxentry{TRAIN}\spxextra{dataset.Split attribute}}

\begin{fulllineitems}
\phantomsection\label{\detokenize{dataset:dataset.Split.TRAIN}}
\pysigstartsignatures
\pysigline
{\sphinxbfcode{\sphinxupquote{TRAIN}}\sphinxbfcode{\sphinxupquote{\DUrole{w}{ }\DUrole{p}{=}\DUrole{w}{ }0}}}
\pysigstopsignatures
\end{fulllineitems}

\index{VAL (dataset.Split attribute)@\spxentry{VAL}\spxextra{dataset.Split attribute}}

\begin{fulllineitems}
\phantomsection\label{\detokenize{dataset:dataset.Split.VAL}}
\pysigstartsignatures
\pysigline
{\sphinxbfcode{\sphinxupquote{VAL}}\sphinxbfcode{\sphinxupquote{\DUrole{w}{ }\DUrole{p}{=}\DUrole{w}{ }1}}}
\pysigstopsignatures
\end{fulllineitems}


\end{fulllineitems}

\index{get\_path() (in module dataset)@\spxentry{get\_path()}\spxextra{in module dataset}}

\begin{fulllineitems}
\phantomsection\label{\detokenize{dataset:dataset.get_path}}
\pysigstartsignatures
\pysiglinewithargsret
{\sphinxcode{\sphinxupquote{dataset.}}\sphinxbfcode{\sphinxupquote{get\_path}}}
{\sphinxparam{\DUrole{n}{pid}}\sphinxparamcomma \sphinxparam{\DUrole{n}{data\_dir}}}
{}
\pysigstopsignatures
\end{fulllineitems}


\sphinxstepscope


\section{hirescam}
\label{\detokenize{hirescam:module-hirescam}}\label{\detokenize{hirescam:hirescam}}\label{\detokenize{hirescam::doc}}\index{module@\spxentry{module}!hirescam@\spxentry{hirescam}}\index{hirescam@\spxentry{hirescam}!module@\spxentry{module}}\index{HiResCam (class in hirescam)@\spxentry{HiResCam}\spxextra{class in hirescam}}

\begin{fulllineitems}
\phantomsection\label{\detokenize{hirescam:hirescam.HiResCam}}
\pysigstartsignatures
\pysiglinewithargsret
{\sphinxbfcode{\sphinxupquote{\DUrole{k}{class}\DUrole{w}{ }}}\sphinxcode{\sphinxupquote{hirescam.}}\sphinxbfcode{\sphinxupquote{HiResCam}}}
{\sphinxparam{\DUrole{n}{model}}\sphinxparamcomma \sphinxparam{\DUrole{n}{device}}\sphinxparamcomma \sphinxparam{\DUrole{n}{model\_name}}\sphinxparamcomma \sphinxparam{\DUrole{n}{target\_layer\_name}}}
{}
\pysigstopsignatures
\sphinxAtStartPar
Bases: \sphinxcode{\sphinxupquote{object}}

\sphinxAtStartPar
Inference layer adapted from \sphinxurl{https://github.com/rachellea/hirescam}, Draelos \& Carin, 2021 (arXiv:2011.08891)
Context: CNNs compute gradients \(K \in \mathbb{R}^{m \times m}\)

\sphinxAtStartPar
Methodology: HiResCAM uses gradients computed by the CNN to produce a heatmap \(L^c_{HiResCAM}
\in \mathbb{R}^{d \times h \times w}\) that highlights important regions in the input volume for a given class \(c\).

\sphinxAtStartPar
See test.py for example usage.
\index{hirescam() (hirescam.HiResCam method)@\spxentry{hirescam()}\spxextra{hirescam.HiResCam method}}

\begin{fulllineitems}
\phantomsection\label{\detokenize{hirescam:hirescam.HiResCam.hirescam}}
\pysigstartsignatures
\pysiglinewithargsret
{\sphinxbfcode{\sphinxupquote{hirescam}}}
{\sphinxparam{\DUrole{n}{target\_grads}\DUrole{p}{:}\DUrole{w}{ }\DUrole{n}{ndarray}}\sphinxparamcomma \sphinxparam{\DUrole{n}{target\_activs}\DUrole{p}{:}\DUrole{w}{ }\DUrole{n}{ndarray}}}
{{ $\rightarrow$ ndarray}}
\pysigstopsignatures\begin{description}
\sphinxlineitem{Compute the HiResCAM heatmap using the formula:}
\sphinxAtStartPar
\(L^c_{\text{HiResCAM}} = \sum_k \frac{\partial Y^c}{\partial A^k} \odot A^k\)

\end{description}
\begin{quote}\begin{description}
\sphinxlineitem{Parameters}\begin{itemize}
\item {} 
\sphinxAtStartPar
\sphinxstyleliteralstrong{\sphinxupquote{target\_grads}} \textendash{} Gradients at the target layer, shape (batch, num\_features, D’, H’, W’)

\item {} 
\sphinxAtStartPar
\sphinxstyleliteralstrong{\sphinxupquote{target\_activs}} \textendash{} Activations at the target layer, shape (batch, num\_features, D’, H’, W’)

\end{itemize}

\sphinxlineitem{Returns}
\sphinxAtStartPar
np.ndarray: Heatmap of shape (batch, D, H, W). NOTE that heatmap renderer is left to the caller.

\end{description}\end{quote}

\end{fulllineitems}

\index{return\_explanation() (hirescam.HiResCam method)@\spxentry{return\_explanation()}\spxextra{hirescam.HiResCam method}}

\begin{fulllineitems}
\phantomsection\label{\detokenize{hirescam:hirescam.HiResCam.return_explanation}}
\pysigstartsignatures
\pysiglinewithargsret
{\sphinxbfcode{\sphinxupquote{return\_explanation}}}
{\sphinxparam{\DUrole{n}{ctvol}}\sphinxparamcomma \sphinxparam{\DUrole{n}{chosen\_label\_index}}}
{}
\pysigstopsignatures
\sphinxAtStartPar
Compute the HiResCam heatmap for a single input volume.
\begin{quote}\begin{description}
\sphinxlineitem{Parameters}\begin{itemize}
\item {} 
\sphinxAtStartPar
\sphinxstyleliteralstrong{\sphinxupquote{ctvol}} \textendash{} Input volume of the shape e.g (1, 3, 32, 256, 256). The caller is responsible for ensuring this is the case.

\item {} 
\sphinxAtStartPar
\sphinxstyleliteralstrong{\sphinxupquote{chosen\_label\_index}} \textendash{} Index of the class for which to compute the heatmap. For binary classification, this will always be 0.

\end{itemize}

\sphinxlineitem{Returns}
\sphinxAtStartPar
Raw CAM volume of shape (1,D,H,W), where D, H, W are the spatial dimensions of the target layer’s output..

\end{description}\end{quote}

\sphinxAtStartPar
The caller is responsible for upsampling and visualizing this heatmap as desired.

\end{fulllineitems}


\end{fulllineitems}


\sphinxstepscope


\section{visualize}
\label{\detokenize{visualize:module-visualize}}\label{\detokenize{visualize:visualize}}\label{\detokenize{visualize::doc}}\index{module@\spxentry{module}!visualize@\spxentry{visualize}}\index{visualize@\spxentry{visualize}!module@\spxentry{module}}
\sphinxAtStartPar
Visualize HiResCAM attention maps overlaid on MRI slices for a single patient.
\begin{description}
\sphinxlineitem{Pipeline:}\begin{enumerate}
\sphinxsetlistlabels{\arabic}{enumi}{enumii}{}{.}%
\item {} 
\sphinxAtStartPar
Load model + one test sample

\item {} 
\sphinxAtStartPar
Run HiResCAM (via hirescam.py) at the last conv layer of the encoder

\item {} 
\sphinxAtStartPar
Upsample the 3D heatmap to the input volume size

\item {} 
\sphinxAtStartPar
For each of the 3 time points, show MRI slices side\sphinxhyphen{}by\sphinxhyphen{}side with the heatmap overlay

\end{enumerate}

\end{description}

\sphinxAtStartPar
Note: The model ends in two FC layers, so HiResCAM’s faithfulness guarantee
(paper Section 3.5) does not strictly apply, but it is still more faithful than Grad\sphinxhyphen{}CAM.
\index{load\_model() (in module visualize)@\spxentry{load\_model()}\spxextra{in module visualize}}

\begin{fulllineitems}
\phantomsection\label{\detokenize{visualize:visualize.load_model}}
\pysigstartsignatures
\pysiglinewithargsret
{\sphinxcode{\sphinxupquote{visualize.}}\sphinxbfcode{\sphinxupquote{load\_model}}}
{\sphinxparam{\DUrole{n}{path}}}
{}
\pysigstopsignatures
\end{fulllineitems}

\index{main() (in module visualize)@\spxentry{main()}\spxextra{in module visualize}}

\begin{fulllineitems}
\phantomsection\label{\detokenize{visualize:visualize.main}}
\pysigstartsignatures
\pysiglinewithargsret
{\sphinxcode{\sphinxupquote{visualize.}}\sphinxbfcode{\sphinxupquote{main}}}
{}
{}
\pysigstopsignatures
\end{fulllineitems}

\index{overlay\_heatmap\_on\_slice() (in module visualize)@\spxentry{overlay\_heatmap\_on\_slice()}\spxextra{in module visualize}}

\begin{fulllineitems}
\phantomsection\label{\detokenize{visualize:visualize.overlay_heatmap_on_slice}}
\pysigstartsignatures
\pysiglinewithargsret
{\sphinxcode{\sphinxupquote{visualize.}}\sphinxbfcode{\sphinxupquote{overlay\_heatmap\_on\_slice}}}
{\sphinxparam{\DUrole{n}{mri\_slice}\DUrole{p}{:}\DUrole{w}{ }\DUrole{n}{ndarray}}\sphinxparamcomma \sphinxparam{\DUrole{n}{heat\_slice}\DUrole{p}{:}\DUrole{w}{ }\DUrole{n}{ndarray}}\sphinxparamcomma \sphinxparam{\DUrole{n}{alpha}\DUrole{o}{=}\DUrole{default_value}{0.5}}}
{}
\pysigstopsignatures
\sphinxAtStartPar
Blend a grayscale MRI slice with a jet\sphinxhyphen{}colormap heatmap.
Both inputs should be in {[}0, 1{]}.
Returns an RGB array.

\end{fulllineitems}

\index{upsample\_heatmap() (in module visualize)@\spxentry{upsample\_heatmap()}\spxextra{in module visualize}}

\begin{fulllineitems}
\phantomsection\label{\detokenize{visualize:visualize.upsample_heatmap}}
\pysigstartsignatures
\pysiglinewithargsret
{\sphinxcode{\sphinxupquote{visualize.}}\sphinxbfcode{\sphinxupquote{upsample\_heatmap}}}
{\sphinxparam{\DUrole{n}{heatmap}\DUrole{p}{:}\DUrole{w}{ }\DUrole{n}{Tensor}}\sphinxparamcomma \sphinxparam{\DUrole{n}{target\_size}\DUrole{p}{:}\DUrole{w}{ }\DUrole{n}{tuple}}}
{{ $\rightarrow$ Tensor}}
\pysigstopsignatures
\sphinxAtStartPar
ReLU + normalize to {[}0,1{]} + trilinear upsample to target\_size (D, H, W).
heatmap: shape (D’, H’, W’)

\end{fulllineitems}


\sphinxstepscope


\section{pay\_attn.model\_outputs}
\label{\detokenize{pay_attn_model_outputs:module-model_outputs}}\label{\detokenize{pay_attn_model_outputs:pay-attn-model-outputs}}\label{\detokenize{pay_attn_model_outputs::doc}}\index{module@\spxentry{module}!model\_outputs@\spxentry{model\_outputs}}\index{model\_outputs@\spxentry{model\_outputs}!module@\spxentry{module}}\index{ModelOutputsPcrCNN (class in model\_outputs)@\spxentry{ModelOutputsPcrCNN}\spxextra{class in model\_outputs}}

\begin{fulllineitems}
\phantomsection\label{\detokenize{pay_attn_model_outputs:model_outputs.ModelOutputsPcrCNN}}
\pysigstartsignatures
\pysiglinewithargsret
{\sphinxbfcode{\sphinxupquote{\DUrole{k}{class}\DUrole{w}{ }}}\sphinxcode{\sphinxupquote{model\_outputs.}}\sphinxbfcode{\sphinxupquote{ModelOutputsPcrCNN}}}
{\sphinxparam{\DUrole{n}{model}\DUrole{p}{:}\DUrole{w}{ }\DUrole{n}{{\hyperref[\detokenize{train:train.PcrCNN}]{\sphinxcrossref{PcrCNN}}}}}\sphinxparamcomma \sphinxparam{\DUrole{n}{target\_layer\_name}\DUrole{p}{:}\DUrole{w}{ }\DUrole{n}{str}}}
{}
\pysigstopsignatures
\sphinxAtStartPar
Bases: \sphinxcode{\sphinxupquote{object}}

\sphinxAtStartPar
Class for running a PcrCNN \textless{}model\textgreater{} and:
(1) Extracting activations from intermediate target layers
(2) Extracting gradients from intermediate target layers
(3) Returning the final model output
\index{get\_gradients() (model\_outputs.ModelOutputsPcrCNN method)@\spxentry{get\_gradients()}\spxextra{model\_outputs.ModelOutputsPcrCNN method}}

\begin{fulllineitems}
\phantomsection\label{\detokenize{pay_attn_model_outputs:model_outputs.ModelOutputsPcrCNN.get_gradients}}
\pysigstartsignatures
\pysiglinewithargsret
{\sphinxbfcode{\sphinxupquote{get\_gradients}}}
{}
{}
\pysigstopsignatures
\end{fulllineitems}

\index{run\_model() (model\_outputs.ModelOutputsPcrCNN method)@\spxentry{run\_model()}\spxextra{model\_outputs.ModelOutputsPcrCNN method}}

\begin{fulllineitems}
\phantomsection\label{\detokenize{pay_attn_model_outputs:model_outputs.ModelOutputsPcrCNN.run_model}}
\pysigstartsignatures
\pysiglinewithargsret
{\sphinxbfcode{\sphinxupquote{run\_model}}}
{\sphinxparam{\DUrole{n}{x}}}
{}
\pysigstopsignatures
\sphinxAtStartPar
\textless{}p\textgreater{}
Runs the model on input x, while saving activations and gradients from the target layer.
\textless{}/p\textgreater{}
:param x: Input x is expected for (batch, 3, 32, 256, 256). The caller is responsible for ensuring this is the case.
:return:

\end{fulllineitems}

\index{save\_gradient() (model\_outputs.ModelOutputsPcrCNN method)@\spxentry{save\_gradient()}\spxextra{model\_outputs.ModelOutputsPcrCNN method}}

\begin{fulllineitems}
\phantomsection\label{\detokenize{pay_attn_model_outputs:model_outputs.ModelOutputsPcrCNN.save_gradient}}
\pysigstartsignatures
\pysiglinewithargsret
{\sphinxbfcode{\sphinxupquote{save\_gradient}}}
{\sphinxparam{\DUrole{n}{grad}}}
{}
\pysigstopsignatures
\end{fulllineitems}


\end{fulllineitems}

\index{return\_model\_outputs\_class() (in module model\_outputs)@\spxentry{return\_model\_outputs\_class()}\spxextra{in module model\_outputs}}

\begin{fulllineitems}
\phantomsection\label{\detokenize{pay_attn_model_outputs:model_outputs.return_model_outputs_class}}
\pysigstartsignatures
\pysiglinewithargsret
{\sphinxcode{\sphinxupquote{model\_outputs.}}\sphinxbfcode{\sphinxupquote{return\_model\_outputs\_class}}}
{\sphinxparam{\DUrole{n}{model\_name}}}
{}
\pysigstopsignatures
\end{fulllineitems}



\renewcommand{\indexname}{Python Module Index}
\begin{sphinxtheindex}
\let\bigletter\sphinxstyleindexlettergroup
\bigletter{d}
\item\relax\sphinxstyleindexentry{dataset}\sphinxstyleindexpageref{dataset:\detokenize{module-dataset}}
\indexspace
\bigletter{h}
\item\relax\sphinxstyleindexentry{hirescam}\sphinxstyleindexpageref{hirescam:\detokenize{module-hirescam}}
\indexspace
\bigletter{m}
\item\relax\sphinxstyleindexentry{model\_outputs}\sphinxstyleindexpageref{pay_attn_model_outputs:\detokenize{module-model_outputs}}
\indexspace
\bigletter{t}
\item\relax\sphinxstyleindexentry{train}\sphinxstyleindexpageref{train:\detokenize{module-train}}
\indexspace
\bigletter{v}
\item\relax\sphinxstyleindexentry{visualize}\sphinxstyleindexpageref{visualize:\detokenize{module-visualize}}
\end{sphinxtheindex}

\renewcommand{\indexname}{Index}
\printindex
\end{document}